\section{The first reflection-free sixfold source degrees}
\label{sec:bireflection-gate}

The statements here concern genus-zero finite-cover sources in the
representation setting of O$''$.  Let $K$ be imaginary quadratic,
let $V$ be an absolutely irreducible $K$-representation of dimension
$m$, and let its complex character field be exactly $K$.  Its finite
image is denoted by $G$.  We consider a generating product-one tuple
$(g_1,\ldots,g_r)$ whose multiplicity variation has $K$-dimension six
and signature $(3,3)$.  Trivial inertia in this representation is
discarded.  We assume that the finite image is reflection-free.
Write $d_j=\operatorname{rank}(1-g_j)$.  The necessary conditions for
dominance are
\[
 d_j\ge2,\qquad \sum_jd_j=2m+6,\qquad r\ge12.
\tag{BI.1}
\]
The last condition is only necessary; none of the arguments below
infers dominance from it.  The exact-character-field hypothesis is
the ordinary absolutely irreducible coefficient setting.  A
rational-character coefficient endowed subsequently with an extra
$K$-action is a separate construction; generic endomorphisms already
obstruct it as in D-E8-GRR-06.

\begin{theorem}[D-BI-01: two-dimensional inducing blocks]
\label{D-BI-01}
Suppose the complex representation is induced from an absolutely
irreducible representation $\tau$ of dimension two of a subgroup
$H$, so $V_{\mathbf C}\simeq\operatorname{Ind}_H^G\tau$.
Then the period map of the associated Weil sixfolds has rank at most
seven.  This assertion does not require the character field of
$\tau$ to equal $K$.
\end{theorem}
\begin{proof}
Let $C\to\mathbf P^1$ be the connected $G$-cover, set $X=C/H$, and
let $H_0$ be the kernel of $\tau$.  The intermediate map
\[
             Y=C/H_0\longrightarrow X
\]
is a connected finite Galois cover with group $H/H_0$.  Frobenius
reciprocity and quotient pullback give an isomorphism of complex
Hodge variations
\[
 \operatorname{Hom}_G(V_{\mathbf C},H^1(C,\mathbf C))
 \simeq \operatorname{Hom}_H(\tau,H^1(Y,\mathbf C)).
\tag{BI.2}
\]
This identification holds in families after a finite cover of the
parameter space marking the subgroup and representations.  It is
induced by fixed finite-group algebra correspondences; it is not
merely an equality of fibre dimensions.

Put $h=g(X)$.  Discard precisely the points whose local inertia is
trivial on $\tau$, and denote the remaining number of branch points
by $R$.  If $D$ is the sum of the moving dimensions of these local
operators on $\tau$, the Euler characteristic of the middle-extension
local system gives
\[
              6=2(2h-2)+D,\qquad D=10-4h.
\tag{BI.3}
\]
The zero-th and second cohomology terms vanish: $\tau$ is a nontrivial
absolutely irreducible representation and the cover has its full
image $H/H_0$.  Faithfulness of this image gives $R\le D$.
In particular $h\le2$.  For $h=0$ the six-dimensional parabolic
cohomology also implies $R\ge5$, since each moving dimension is
at most two.  For $h=1$ one has $R\ge3$.  Thus all pointed curves
in question are stable; the unpointed exceptional moduli dimensions
do not occur.

With the finite group and topological branching data fixed, the
Hurwitz space of $Y\to X$ has finite fibres over the moduli space of
the $R$-pointed genus-$h$ curve: the punctured surface group is
finitely generated and has only finitely many homomorphisms into
the fixed finite group.  Locally this is the usual finite cover of
the pointed Teichmuller space.  Consequently (BI.2) factors through
a space of dimension at most
\[
       3h-3+R\le3h-3+(10-4h)=7-h\le7.
\tag{BI.4}
\]
Forgetting the additional map $X\to\mathbf P^1$ can lose parameters,
but cannot add variation to (BI.2).

Finally, the one complex multiplicity space in (BI.2) determines the
period of the rational $K$-Hodge structure.  Indeed, its filtration
$F^1_\sigma$ determines $F^1_{\bar\sigma}$ as its annihilator under
the fixed polarization pairing between the two embeddings of $K$.
Thus the bound is a bound for the full Weil period map, even when
$\tau$ itself is defined over a larger field.  No descent of
$\tau$ to $K$ was used.
\end{proof}

\begin{lemma}[D-BI-02a: tensor products in degrees nine through fourteen]
\label{D-BI-02a}
Under (BI.1), if $9\le m\le14$, the representation is tensor
indecomposable in the sense required by Guralnick--Saxl's
Theorem~7.1.
\end{lemma}
\begin{proof}
Since $2m+6<36$, (BI.1) forces at least one $d_j=2$.
Suppose there is a factorwise invariant tensor decomposition
$V_{\mathbf C}=A\otimes B$, with $2\le a=\dim A\le b=\dim B$.
The factors of any finite-order tensor operator can be rescaled to
be diagonalizable of finite order.  If one factor is nonscalar,
its largest eigenspace has codimension at least one.  Summing
products of eigenvalue multiplicities shows that a nonscalar tensor
operator has moving dimension at least $\min(a,b)$, even after
an arbitrary scalar multiplication.  Thus the existence of a
bireflection excludes $a\ge3$ in the stated degree range.

It remains that $a=2$ and $b\in\{5,6,7\}$.  The projective image
on $A$ is irreducible, since an invariant line in $A$ would give an
invariant proper subspace of $A\otimes B$.  At least three of the
inertia operators have nontrivial image in $\operatorname{PGL}(A)$.
Otherwise their product-one relation makes that image trivial or
cyclic, and a finite cyclic projective subgroup of
$\operatorname{PGL}_2(\mathbf C)$ fixes a line.  Each of these
three operators has moving dimension at least $b$: both eigenvalue
multiplicities on $A$ are one, so its tensor product has a fixed
space of dimension at most $b$.  The other inertia operators have
moving dimension at least two.  Hence
\[
       2m+6\ge3b+2(r-3)\ge\frac{3m}{2}+18,
\]
which requires $m\ge24$, a contradiction.

The word ``tensor indecomposable'' in the cited theorem excludes
factorwise tensor decompositions.  Groups permuting equal tensor
factors, for example $\operatorname{GL}(A)\wr\mathfrak S_t$ on
$A^{\otimes t}$, are explicitly included in the theorem's
tensor-indecomposable case; see its introductory paragraph on
p.~548.  Thus no tensor-permutation case has been omitted here.
\end{proof}

\begin{theorem}[D-BI-02: primitive degrees nine through fourteen]
\label{D-BI-02}
There is no source satisfying (BI.1) with primitive complex
representation and $9\le m\le14$ whose character field is imaginary
quadratic.  This exclusion does not depend on its determinant class.
\end{theorem}
\begin{proof}
There is a bireflection and D-BI-02a applies.  We use exactly the
characteristic-zero specialization of the published
Guralnick--Saxl classification \cite[Theorem~7.1 and Remark~(3),
pp.~548--549]{BI:GS}: in these dimensions the possibilities are
scalar extensions of the deleted permutation representation of an
alternating or symmetric group, or the ten-dimensional representation
of $U_5(2)$ and its linear normalizer.  The defining-characteristic
classical cases in that theorem do not occur over $\mathbf C$.

In the permutation case put $n=m+1$.  After conjugation every group
element is $\lambda\sigma$ on the deleted subspace of
$\mathbf C^n$, with $\sigma\in\mathfrak S_n$ and $\lambda$ a root
of unity.  This description includes the normalizer of
$\mathfrak A_n$: for $n\ge10$ its automorphisms are induced by
$\mathfrak S_n$, and Schur's lemma supplies only an extra scalar.
If $\lambda\notin\{1,-1\}$ has order $e\ge3$, the
$\lambda^{-1}$-eigenspace of a permutation matrix has dimension
at most $\lfloor n/e\rfloor\le\lfloor n/3\rfloor$.
Indeed a cycle contributes at most one such eigenvector and must
have length divisible by $e$.  Deleting the trivial line does not
change this nontrivial eigenspace.  Thus
\[
 \operatorname{rank}(1-\lambda\sigma)
       \ge q_m:=m-\left\lfloor\frac{m+1}{3}\right\rfloor.
\tag{BI.5}
\]
If all inertia scalars are real, their generated representation is
real and its character, lying in $K$, is rational.  Otherwise at
least two inertia scalars are nonreal: the product of the scalars
is one, since the deleted permutation group contains no nontrivial
scalar.  Equation (BI.1) would then imply
\[
        2m+6\ge2q_m+2(r-2)\ge20+2q_m.
\tag{BI.6}
\]
The following exact values contradict this inequality:
\[
\begin{array}{c|rrrrrr}
m&9&10&11&12&13&14\\ \hline
2m+6&24&26&28&30&32&34\\
q_m&6&7&7&8&9&9\\
20+2q_m&32&34&34&36&38&38.
\end{array}
\]

For $U_5(2)$ one has $m=10$.  Its degree-ten module preserves a
nondegenerate alternating form.  This is the symplectic Weil module
in \cite[Table~II, author version p.~7; Proposition~7.10]{BI:GT}; it is also
given by explicit matrices and an alternating form of determinant
one in the author's ATLAS data \cite{BI:AtlasU52}.
Irreducibility makes the invariant form line unique, so every
element of its linear normalizer is a similitude of this form.
If the multiplier of an inertia operator is not one, its fixed
space is totally isotropic: for fixed vectors $x,y$ one has
$b(x,y)=b(gx,gy)=c(g)b(x,y)$.
Its fixed dimension is therefore at most five, and its moving
dimension is at least five.  But (BI.1) bounds every moving
dimension by
\[
              26-2(r-1)\le4.
\]
Every inertia multiplier is consequently one.  They generate $G$,
so $G$ preserves the alternating form, the representation is
self-dual, and its character is real.  Once again an imaginary
quadratic character field is impossible.
\end{proof}

\begin{proposition}[D-BI-03: the full degree-nine nonsplit exclusion]
\label{D-BI-03}
No reflection-free source of coefficient degree at most nine can
dominate a nonsplit Weil sixfold family in this representation
setting.  Thus the least numerical possibility $m=9$ in (BI.1)
does not occur.
\end{proposition}
\begin{proof}
If $m<9$, then $2m+6<24$, contrary to (BI.1).  If $m=9$, then
necessarily $r=12$ and every inertia has moving dimension two.
The primitive case is D-BI-02.  An imprimitive irreducible
nine-dimensional representation has either nine one-dimensional
blocks or three three-dimensional blocks.  The first case is
monomial and is excluded for a dominant nonsplit sixfold by
D-MON-01--03.  In the second case, an operator nontrivially
permuting the blocks has moving dimension at least three: a
cycle of length $\ell\ge2$ on blocks of dimension three has
fixed dimension at most three and moving dimension at least
$3(\ell-1)$.  Every inertia therefore fixes every block.
Since the tuple generates $G$, this contradicts irreducibility.
\end{proof}

\begin{proposition}[D-BI-04: the remaining imprimitive blocks below fifteen]
\label{D-BI-04}
Suppose (BI.1) holds with $9\le m\le14$, and suppose the source
could dominate a nonsplit sixfold family.  Apart from the excluded
monomial and two-dimensional-block cases, the only imprimitive
block sizes still permitted are
\[
            (m,t,d)=(12,4,3)\quad\hbox{or}\quad(14,2,7),
\tag{BI.7}
\]
where $m=td$.  The intermediate block quotient $X\to\mathbf P^1$
has genus zero.  Equality holds in every rank estimate: the first
case has six block-transposition inertia operators of rank three
and six block-fixed inertia operators of rank two; the second has
two block-transposition operators of rank seven and ten block-fixed
operators of rank two.  This proposition is a reduction, not a
claim that either remaining case exists or dominates.
\end{proposition}
\begin{proof}
For a transitive imprimitivity system with $t$ blocks of dimension
$d$, let $\pi_j$ be the permutation of the blocks induced by
$g_j$, and let $\operatorname{ind}(\pi_j)=t-\#\text{cycles}(\pi_j)$.
The fixed space on a block cycle has dimension at most $d$, so
\[
             d_j\ge d\operatorname{ind}(\pi_j).
\]
If $\pi_j\ne1$, then
$d_j-2\ge(d-2)\operatorname{ind}(\pi_j)$; the difference in this
estimate includes $2(\operatorname{ind}(\pi_j)-1)$.
For $\pi_j=1$ the same inequality follows from $d_j\ge2$.
If $h=g(X)$ for the block quotient, Riemann--Hurwitz gives
\[
 \sum_j\operatorname{ind}(\pi_j)=2t+2h-2.
\]
Consequently
\[
 2m-18\ge\sum_j(d_j-2)
       \ge(d-2)(2t+2h-2),
\qquad
 2t+d\ge11+(d-2)h.
\tag{BI.8}
\]
For $d\ge3$, $t\ge2$ and $9\le td\le14$, the last inequality
leaves exactly the two pairs in (BI.7).  Both satisfy $2t+d=11$,
so $h=0$ and every inequality used is an equality.  In particular
$r=12$, every nontrivial block permutation has index one and hence
is a transposition, each has moving rank exactly $d$, and every
block-fixed inertia has moving rank two.  The required number of
transpositions is $2t-2$, which gives the two stated patterns.
\end{proof}

\begin{theorem}[D-BI-05: three-dimensional inducing blocks]
\label{D-BI-05}
Let $K$ be one of
\[
 \mathbf Q(i),\quad \mathbf Q(\sqrt{-2}),\quad
 \mathbf Q(\sqrt{-7}),\quad\mathbf Q(\sqrt{-11}),\quad
 \mathbf Q(\sqrt{-19}),\quad\mathbf Q(\sqrt{-43}),\quad
 \mathbf Q(\sqrt{-67}),\quad\mathbf Q(\sqrt{-163}).
\]
In the exact-character-field setting above, an irreducible
representation induced from a three-dimensional complex representation
cannot dominate a nonsplit Weil sixfold family.  The inducing
representation need not have a model over $K$, and the original
coefficient degree is unrestricted.
\end{theorem}
\begin{proof}
Write $V_{\mathbf C}=\operatorname{Ind}_H^G\tau$ with
$\dim\tau=3$ and $L=\mathbf Q(\chi_\tau)$.  Induction of
characters immediately gives
\[
                         K=\mathbf Q(\chi_V)\subset L.
\tag{BI.9}
\]
This inclusion uses values of complex characters, so does not
presuppose a rational block system.  Form the intermediate cover
$Y=C/\ker\tau\to X=C/H$ as in D-BI-01.  Its multiplicity variation
is the same one complex period variation as the original family.
If $h=g(X)$, $R$ is its genuine branch count, and $D$ the sum of
its moving dimensions, Euler characteristic gives
\[
                   D=12-6h,\qquad R\le D.
\]
The stable pointed-curve Hurwitz dimension is at most
\[
                    3h-3+R\le9-3h.
\tag{BI.10}
\]
The potentially unstable cases do not give an exception: at genus
zero the positive six-dimensional parabolic cohomology forces at
least four punctures; at genus one $D=6$ forces branching; at genus
two the unramified base has dimension three.  Thus rank nine forces
$h=0$, $R=12$, and equality $R=D$.  Every inertia on $\tau$ is
therefore a complex reflection.  Since $X=\mathbf P^1$, these
inertia generate its full image, which is a finite irreducible
rank-three reflection group.

If $\tau$ is imprimitive, its rank being three makes it monomial;
induction then makes $V$ monomial.  D-MON-01--03 excludes a
dominant nonsplit source in this case.  Otherwise the published
rank-three reflection classification and its splitting fields give
\[
\begin{array}{c|ccccc}
\text{group}&G_{23}=W(H_3)&G_{24}&G_{25}&G_{26}&G_{27}\\ \hline
k&\mathbf Q(\sqrt5)&\mathbf Q(\sqrt{-7})&
 \mathbf Q(\omega)&\mathbf Q(\omega)&
 \mathbf Q(\omega,\sqrt5).
\end{array}
\tag{BI.11}
\]
Here the classification and the $G_{25}$ realization are
\cite[\S4.16, p.~412; Theorem~5.12, p.~435]{FIELD:Cohen};
the splitting fields are \cite[Table~8.2, p.~675, with its
omission convention on p.~676]{FIELD:Malle}.  In each case
$L\subset k$, which is all that is needed.  The quadratic
subfields of the last biquadratic field are
$\mathbf Q(\sqrt5)$, $\mathbf Q(\sqrt{-3})$, and
$\mathbf Q(\sqrt{-15})$.  Hence for the eight fields in the
statement, (BI.9) leaves only $G_{24}$ and
$K=L=\mathbf Q(\sqrt{-7})$.

In this last case the natural $G_{24}$ representation has a model
over $K$.  Both $V$ and $\operatorname{Ind}_H^G\tau$ then have
$K$-models and equal absolutely irreducible complex characters;
their intertwining linear equations over $K$ give a $K$-isomorphism.
Shapiro's identification of the parabolic complexes consequently
identifies their polarized $K$-variations up to a positive rational
scalar: use the induced positive Hermitian form and sum cup products
over cosets; the specified irreducible invariant form differs by
such a scalar by Schur's lemma.  This is the polarization argument
of D-FIELD-03, now justified after establishing $L=K$.
On a six-dimensional $K$-space a rational scalar changes the
determinant by its sixth power, a $K$-norm.  D-DET-02 proves that
the natural $G_{24}$ source in question has split signed determinant.
The original source is therefore split as well, a contradiction.
\end{proof}

\begin{proposition}[D-BI-06: the remaining two blocks of dimension seven]
\label{D-BI-06}
The $(m,t,d)=(14,2,7)$ case left by D-BI-04 cannot dominate a
nonsplit Weil sixfold family with exact imaginary quadratic
character field.
\end{proposition}
\begin{proof}
Let $\tau$ be the representation on one block.  If it is
imprimitive, its prime degree seven makes it monomial, hence the
original representation is monomial; the nonsplit case is already
excluded.  We can thus assume $\tau$ is primitive.  D-BI-04 gives
$X=\mathbf P^1$, two minimal-rank block transpositions, and ten
block-fixed rank-two inertia operators.  A minimal-rank block
transposition has the form
\[
       \begin{pmatrix}0&A\\A^{-1}&0\end{pmatrix},
\]
because its fixed space has dimension seven.  Its square is the
identity, so it induces no inertia on $\tau$ above the corresponding
ramification point of $X\to\mathbf P^1$.  Every remaining inertia
on $\tau$ comes from a block-fixed operator and has rank at most two.

The intermediate six-dimensional multiplicity has total moving
dimension $2\cdot7+6=20$.  Rank nine of its period map requires at
least twelve genuine branch points on $X$.  Thus at least one
inertia on $\tau$ is a reflection; otherwise the branch count is
at most ten.  Let $R_0$ be the normal closure of all reflections
in the finite image of $\tau$.  Primitivity makes $R_0$ irreducible:
its fixed space is invariant, and its orthogonal reflection
components would otherwise be a nontrivial imprimitivity system.

By the complex reflection classification in rank seven, $R_0$ is
real or belongs to the imprimitive series.  An imprimitive group
in that rank has a unique imprimitivity system, by
\cite[Lemma~2.7, pp.~389--390]{FIELD:Cohen}; none of the four
exceptions there has rank seven.  Its normalizer would preserve
that system, contrary to primitivity of $\tau$.  Therefore $R_0$
has a real reflection representation.  The same conclusion covers
the standard $A_7$ representation as well as the exceptional $E_7$
representation; neither is being omitted.

Choose its real positive definite form $B$.  Uniqueness of the
invariant Hermitian form for the irreducible $R_0$ shows that the
finite image of $\tau$ is unitary for $B$.  If $T$ belongs to
this finite image, then $T^{-1}\bar T$ centralizes $R_0$, since the automorphism
of $R_0$ induced by $T$ has real matrices.  Schur's lemma writes
$T=\lambda A$ with $A$ real; rescaling by a positive real number
makes $A$ orthogonal for $B$ and $|\lambda|=1$.
If $\lambda$ is nonreal, its inverse is a nonreal eigenvalue of
the real seven-dimensional $A$.  Conjugate eigenvalues have the
same multiplicity, so this multiplicity is at most three.  Hence
$\operatorname{rank}(1-T)\ge4$.  Every inertia on $\tau$ has rank
at most two, so it is real.  These inertia generate the full image
over $X=\mathbf P^1$, making $\mathbf Q(\chi_\tau)$ real.
But $\chi_V=\operatorname{Ind}_H^G\chi_\tau$ implies
$K\subset\mathbf Q(\chi_\tau)$, a contradiction.
\end{proof}

\begin{corollary}[D-BI-07: a strictly smaller remaining source gate]
\label{D-BI-07}
For the eight non-Eisenstein class-number-one fields in D-BI-05,
a dominant nonsplit finite-cover source with exact character field
$K$ and reflection-free image must have coefficient degree at least
fifteen.  It cannot be induced from a complex representation of
dimension one, two, or three.  The Eisenstein field is omitted
because its all-determinant sixfold theorem is already established
elsewhere in this paper, not because the field was discarded from
the original target.
\end{corollary}
\begin{proof}
Degrees below nine fail (BI.1).  D-BI-02 excludes the primitive
degrees nine through fourteen.  For the imprimitive degrees in
that range, one-dimensional blocks are D-MON-01--03,
two-dimensional blocks are D-BI-01, and D-BI-04 leaves only
the cases excluded by D-BI-05 and D-BI-06.  Inducing dimensions
one, two and three in arbitrary original degree are excluded by
the same three results.
\end{proof}

\paragraph{Verification and next gate.}
The two-dimensional-block factorization was independently derived
by the imprimitive agent.  The character-field inclusion and its
application to all three-dimensional inducing blocks were independently
pointed out by root.  The other arguments are new deductions from
(BI.1), the stated published finite classification, and elementary
eigenvalue estimates.  The companion exact checker verifies the
permutation and block-budget tables and the invariant forms in the
retained degree-ten ATLAS matrices.  It does not claim to prove the
classification or replace the family-level factorization proof.
Degrees at least fifteen with no inducing representation of dimension
at most three are the remaining reflection-free gate for the eight
fields in D-BI-07.  This is a source obstruction, not a proof of
algebraicity for those remaining determinant families.
